\documentclass[10pt,a4paper]{amsart}
\usepackage[margin=19mm]{geometry}
\usepackage{amsmath,amssymb,mathtools}
\usepackage[scr=rsfs,cal=euler]{mathalfa}
\usepackage{xfrac}
\usepackage[unicode,hidelinks,bookmarks=false]{hyperref}
\newcommand{\ie}{i.e.\ }
\newcommand{\cf}{cf.\ }
\newcommand{\resp}{respectively\ }
\newcommand{\Subsection}{\subsection}
\providecommand{\nobreakdash}{\nobreak\hbox{-}}
\setlength{\emergencystretch}{2em}
\multlinegap0pt
\usepackage{amssymb}
\usepackage{tikz}
\newtheorem{theorem}{Theorem}
\newtheorem{remark}{Remark}
\newtheorem{proposition}{Proposition}
\renewcommand\geq{\geqslant}
\renewcommand\leq{\leqslant}
\title{Freeness of arrangements of lines and one conic with ordinary quasi-homogeneous singularities}
\author{Piotr Pokora}
\date{Source: arXiv:2312.13052v5 (CC BY 4.0)}
\begin{document}
\maketitle

\noindent\emph{Source and licence.} Freeness of arrangements of lines and one conic with ordinary quasi-homogeneous singularities. Version arXiv:2312.13052v5 (CC BY 4.0), distributed by arXiv under the Creative Commons Attribution 4.0 International licence, http://creativecommons.org/licenses/by/4.0/, as recorded in the arXiv licence metadata for this exact version and shown on the abstract page https://arxiv.org/abs/2312.13052v5. Full licence notice: \texttt{licenses/arxiv-2312.13052-CC-BY-4.0.txt}.\par\smallskip

\noindent\textit{Editorial scope.} Counted unit: the article's proposition proposition (source0.tex lines 281--284). The article's complete original proof of it is delivered with it (source0.tex lines 285--290, one \texttt{proof} environment of 6 lines). No mathematics is written, repaired or re-derived: the statement and the proof are the article's own lines, in the article's own notation, and no retained line is substituted or re-flowed.  Retained uncounted as the source-local context the proof needs: lines 245--253; lines 264--267.  Nothing was omitted from any retained span, so the package carries no omission record.  No retained line contains a citation, so the wrapper carries no bibliography and no bibliography entry.  No retained line contains a figure environment, an image-inclusion command or drawing code, so no image is delivered and the package declares no asset dependency.  Editorial formatting only, in addition: an AMS \texttt{amsart} preamble that restates the article's own declarations and macros used by the retained lines, namely the environments \texttt{theorem}, \texttt{remark}, \texttt{proposition} on separate counters, together with the standard packages those lines need.  The retained environments print in this document as Theorem 1, Remark 1, Proposition 1 in order of appearance, whereas the article prints the same items with the numbering of its own full document.  The complete native source of the article as distributed in the arXiv e-print for this exact version is bundled unchanged as \texttt{sources/151829/source0.tex}.
\begin{theorem}
\label{dddp}
Let $C = \{f=0\}$ be a reduced curve in $\mathbb{P}^{2}_{\mathbb{C}}$. One has
\begin{equation}
\label{duPles}
(d-1)^{2} - r(d-r-1) = \tau(C)
\end{equation}
if and only if $C = \{f=0\}$ is a free curve, and then $r \leq (d-1)/2$.
\end{theorem}

\begin{remark}
\label{lct}
If $p=(0,0) \in \mathbb{C}^{2}$ is an ordinary singularity of multiplicity $r$ determined by $C = \{ f=0\}$, then ${\rm lct}_{p}(f) = \frac{2}{r}$.
\end{remark}

\begin{proposition}
Let $\mathcal{CL} = \{f=0\}$ be an arrangement of $d\geq 3$ lines and one smooth conic that admits only ordinary singularities of multiplicity $<5$. Assume that $\mathcal{CL}$ is free, then
$${\rm mdr}(f) \in \bigg\{\bigg\lceil \frac{d-2}{2} \bigg\rceil, \bigg\lfloor \frac{d+1}{2} \bigg \rfloor \bigg\}.$$
\end{proposition}

\begin{proof}
Since $\mathcal{CL}$ is free, then by Theorem \ref{dddp} one has ${\rm mdr}(f) \leq \lfloor \frac{d+1}{2} \rfloor$. Moreover, if $\mathcal{CL}$ admits only ordinary singularities with multiplicities $<5$, then $\alpha_{\mathcal{CL}} = \frac{1}{2}$, and this follows from Remark \ref{lct}. This gives us that
$${\rm mdr}(f) \geq \frac{d+2}{2} - 2 = \frac{d-2}{2},$$
and we finally obtain
$${\rm mdr}(f)\geq \bigg\lceil \frac{d-2}{2} \bigg\rceil.$$
\end{proof}

\end{document}
